% !TeX root = ../main.tex
\chapter{Cumprimento de metas FIFO e matching de doadores}
\label{ap:fulfillment}

Este apêndice documenta o núcleo algorítmico da solução: a alocação \textit{First In, First Out} (FIFO) de doações concluídas às solicitações ativas de um hemocentro e o serviço de domínio que identifica doadores elegíveis para notificação.

\section{Alocação FIFO de doações}

A Listagem~\ref{lst:fifo} sintetiza o caminho de leitura do \codigo{DonationRequestFulfillmentService}: \codigo{fill} recebe o hemocentro e a data de referência \codigo{asOfDate}, carrega somente as solicitações ativas e não expiradas daquele centro, e \codigo{allocateLoaded} busca as doações concluídas na janela \codigo{[pedido ativo mais antigo, asOfDate]}. Em seguida, \codigo{allocateFifo} isola o pool por hemocentro, ordena solicitações e doações da mais antiga para a mais nova e atribui cada doação à primeira solicitação que ainda cabe na meta e que \codigo{acceptsDonation} aceita --- o \texttt{break} encerra a busca, concretizando a política FIFO. O resultado permanece em memória: nada é persistido.

\lstinputlisting[
  language=Java,
  caption={Alocação FIFO em \texttt{DonationRequestFulfillmentService}},
  label={lst:fifo}
]{apendices/codigo/FulfillmentDistributeDonations.java}
\fonte{Elaborado pelos próprios autores (2026).}

\section{Identificação de doadores elegíveis}

A Listagem~\ref{lst:matching} apresenta \codigo{DonorMatchingService.findEligibleDonors}, que compõe compatibilidade tipológica (\codigo{canBeFulfilledBy}), elegibilidade temporal (\codigo{isEligibleToDonate}) e proximidade geográfica limitada pela distância máxima configurada pelo doador.

\lstinputlisting[
  language=Java,
  caption={Método \texttt{DonorMatchingService.findEligibleDonors}},
  label={lst:matching}
]{apendices/codigo/DonorMatchingService.java}
\fonte{Elaborado pelos próprios autores (2026).}
